\hypertarget{excluded-hedge-components-and-rebalancing}{%
\section{Excluded hedge components and rebalancing}\label{excluded-hedge-components-and-rebalancing}}

The following conditional methods use the retained AASB~9 compilation for calendar-2024/2025 Australian for-profit Tier~1 illustrations, without optional later amendments. Actual authority admission and component valuation remain separate requirements.

\hypertarget{the-option-premium-pays-for-protection}{%
\subsection{The option premium pays for protection}\label{the-option-premium-pays-for-protection}}

A purchased option can protect an inventory purchase even if it expires with no intrinsic value. Its time value is the price of that protection. When only intrinsic value is designated, applying the earlier whole-derivative method would mix the price of protection with the payoff that offsets the purchase. AASB~9.6.5.15 \citep[6.5.15 and B6.5.29--33]{aasb9dated} separates them. The designated intrinsic-value movement follows the applicable hedge method; the related time-value movement accumulates in a different equity component. The nature of the hedged item determines when that cost affects earnings. A transaction such as an inventory purchase uses a cost attached to that transaction. Protection over a period uses systematic, rational amortisation over the period in which the intrinsic-value hedge adjustment can affect earnings. The life of the option alone does not decide which route applies (B6.5.29--30).

The event record must identify the instrument, designated component, underlying item and risk, critical terms, actual and aligned time values at designation and every measurement date, intrinsic-value changes, separate reserve balances, the transaction or time-period classification, amortisation policy, and each basis transfer or reclassification. Existing intrinsic-value effectiveness requirements remain necessary. Unknown component valuations cannot be replaced with a total derivative value and a zero time value.

Suppose a qualifying inventory hedge begins with cash and equity of 200. Buy an option for 6, entirely time value. Cash is 194 and the option asset is 6; there is no initial expense. The first supplied valuation is intrinsic value 12 and aligned time value 4. The option is worth 16. Its increase of 10 comprises an effective intrinsic gain of 12 and a time-value loss of 2, both in OCI but in separate reserves. At purchase and expiry, intrinsic value is 20 and time value zero. The option is worth 20, the intrinsic reserve is 20 and the cost reserve is negative 6. The last period's change of 4 comprises intrinsic gain 8 and time-value loss 4.

Buy the inventory for 120 and settle the option for 20. Settlement removes the option asset and leaves cash 94. The intrinsic reserve reduces inventory cost by 20. Removing the negative cost reserve adds 6 to inventory cost:

\[120 - 20 + 6 = 106.\]

Both direct transfers affect the asset and equity; neither is a second cash flow or an OCI reclassification. Cash 94 plus inventory 106 equals equity 200 after the transfers. Selling for 130 gives profit 24 and cash 224. The cash bridge is \(200-6-120+20+130=224\). The equity bridge is opening 200, profit 24, cumulative OCI 14, and a direct basis transfer of negative 14. Omitting the premium would report cost 100 and overstate profit by 6. For a forecast sale instead, the accumulated cost enters earnings when the sale does. An unrecoverable portion enters earnings immediately under 6.5.15(b)(iii), rather than waiting for a transaction that cannot recover it.

\hypertarget{protection-over-time-and-imperfectly-matching-terms}{%
\subsection{Protection over time and imperfectly matching terms}\label{protection-over-time-and-imperfectly-matching-terms}}

Consider separately an option that protects an existing inventory holding for two equal periods. Admit aligned initial time value 6 and a rational even allocation over those periods. Intrinsic-value and inventory movements are zero in this illustration, isolating the protection cost. After one period, supplied time value 2 gives an OCI loss of 4. Amortisation of 3 is reclassified to expense, leaving the cost reserve at negative 1. After expiry, the final time-value loss of 2 and second amortisation of 3 leave no reserve; cumulative expense is 6. Immediately expensing each observed time-value fall would instead charge 4 and 2. That is a different timing policy. If the intrinsic-value hedging relationship is discontinued after the first period, 6.5.15(c) requires the remaining net reserve loss of 1 to enter earnings immediately. The remaining option asset of 2 is still an asset; discontinuation does not destroy it.

When critical terms differ, B6.5.32 requires an aligned option valuation using terms that perfectly match the hedged item. The hypothetical option is a measurement input for the related cost, not a replacement of the actual contract. Under B6.5.33(a) \citep[B6.5.32--33]{aasb9dated}, if actual time value exceeds aligned time value at inception, base the cost reserve on aligned changes and put the difference between actual and aligned changes in profit or loss. For initial actual/aligned values 6/4 and later values 3/3, actual loss 3 separates into OCI loss 1 and profit loss 2. At expiry, cumulative losses 6/4 leave OCI loss 4 and profit loss 2, before the required transfer or reclassification.

If actual time value is initially lower, B6.5.33(b) limits the reserve by the lower cumulative change of actual and aligned time value; the remainder of actual changes enters earnings. In a same-direction decline example, initial values 2/3 become 1/2.5. Their cumulative losses are 1 and 0.5, so the cost reserve holds loss 0.5 and profit holds loss 0.5. At expiry, cumulative losses are 2 and 3: the reserve loss becomes 2 and cumulative profit from the mismatch returns to zero. The second period therefore records OCI loss 1.5 and a profit gain of 0.5, while actual value falls 1. Applying a fresh lower-of rule only to each period would lose this reversal. These numbers test declining positive time values; they do not define a signed-value selector for every collar or complex combination. Such cases must retain the source-defined alignment and valuation analysis rather than infer it from this small example.

A zero-cost collar can have zero initial time value and nonzero intermediate time value. B6.5.31 still requires the related changes in OCI. A supplied movement from zero to 1.2 and back to zero produces OCI gain 1.2 and loss 1.2; zero initial premium is not permission to omit both. The terminal transaction cost adjustment and initial-cost amortisation are zero in that example.

\hypertarget{forward-elements-and-currency-basis-spreads}{%
\subsection{Forward elements and currency basis spreads}\label{forward-elements-and-currency-basis-spreads}}

AASB~9.6.5.16 \citep[6.5.16 and B6.5.34--39]{aasb9dated} permits the corresponding treatment when the forward element is excluded from a spot-designated forward, or a foreign-currency basis spread is excluded. The accounting election is explicit; designation of spot alone does not automatically make the excluded component an OCI reserve. Apply the transaction/time-period distinction and aligned-value rules in B6.5.34--39. For the forward element, the inception comparison uses absolute actual and aligned amounts. Retain their signed fair values and the changes used for the reserve; do not compare signed inception amounts as though a negative forward element were necessarily smaller in economic magnitude. B6.5.39 carries the same guidance to currency basis spreads. The exact valuation of unmatched forward elements and collars remains a required method for those applications.

\hypertarget{rebalancing-preserves-the-history-of-the-hedge}{%
\subsection{Rebalancing preserves the history of the hedge}\label{rebalancing-preserves-the-history-of-the-hedge}}

A changed hedge ratio can preserve the original risk-management objective. It does not erase the accounting before the change. Under 6.5.5--6.5.6 and B6.5.7--21 \citep[6.5.5--6.5.6 and B6.5.7--21]{aasb9dated}, first measure and recognise ineffectiveness under the existing designation, then adjust designated quantities prospectively. Distinguish ordinary fluctuations around a valid ratio from a changed relationship that requires a new ratio. A changed objective requires discontinuation of that relationship; a counterparty default is not repaired by changing quantities. Update the documented expected sources of ineffectiveness after rebalancing.

Before a change, cumulative derivative gain 12 and offsetting item loss 10 give reserve 10 and profit 2. Immediately before rebalancing, admitted values 18 and negative 14 instead give reserve 14 and cumulative profit 4. Record the period's derivative gain 6 as reserve gain 4 and profit gain 2 before adjusting the quantities. A reset to equal amounts at the new designation date would hide the accumulated ineffectiveness.

Store each retained, added and removed volume with its designation date and valuation reference. An added item layer begins at the rebalancing date; the earlier layer keeps its old reference. The source's example has 100 tonnes designated at forward price 80 and adds 10 tonnes when the price is 90. In a zero-discount illustrative comparison at price 95, their changes are \(100(95-80)+10(95-90)=1550\). Applying the old price to all 110 tonnes gives 1650, inventing a gain of 100 on the added layer. Actual measurement must use the admitted risk, dates, discounting and contractual terms.

Adding an instrument similarly starts its change at the new designation date without resetting the old instrument. Decreasing designated instrument volume leaves any retained undesignated derivative at fair value through profit or loss unless another qualifying designation applies. For example, removing 10\% of a homogeneous instrument from the designation leaves its total carrying value unchanged; of a subsequent total fair-value gain 10, the undesignated 1 enters profit and the designated 9 follows the continuing hedge's effectiveness test. This split does not assert that all 9 is effective. Decreasing item volume invokes partial-discontinuation rules for the removed part. Whether its forecast cash flows remain expected determines the reserve's later release; removal from designation alone does not mean cancellation.

These procedures extend the existing supplied-valuation method. They leave delayed or contingent inducement recognition, award conditions and settlement choices, broader EPS participation, group/net-position hedges and complex component valuation as named financing work, rather than treating a balanced illustration as the completed financing design.
